\chapter{KẾT LUẬN VÀ HƯỚNG PHÁT TRIỂN}
\label{chapter:ket_luan}

\section{Kết luận}
Đề tài tốt nghiệp ``Nghiên cứu phương pháp RAG đa phương thức cho trợ lý ảo thông minh hỗ trợ dạy và học môn Khoa học tự nhiên cấp Trung học cơ sở'' đã được triển khai một cách hệ thống và toàn diện: từ thiết lập cơ sở lý luận, thiết kế kiến trúc phân tầng, chuẩn hóa quy trình trích xuất học liệu đa phương thức, đến tiến hành thực nghiệm đối chiếu định lượng và phân tích chuyên sâu. Đối chiếu với \textbf{năm mục tiêu nghiên cứu} được xác lập tại Chương~\ref{chapter:tong_quan_de_tai}, các kết quả trọng tâm đạt được như sau:

\begin{enumerate}
    \item \textbf{Mục tiêu 1 về xây dựng cơ chế truy vấn đa định dạng:} Đạt được phần lớn các yêu cầu đề ra. Hệ thống đã hiện thực hóa thành công khả năng tiếp nhận và truy xuất đồng thời dữ liệu văn bản lý thuyết cùng tài nguyên hình ảnh với $3.881$ khung hình được bóc tách minh bạch theo phương pháp neo chú thích chính quy. Đối với bài toán đặc thù về công thức định lượng thuộc các phân môn Vật lí và Hoá học, đề tài đã nhận diện rõ điểm nghẽn của công cụ nhận dạng ký tự quang học truyền thống, đồng thời đề xuất và triển khai thành công kiến trúc xử lý lai phân cấp theo dòng trên toàn bộ kho tri thức ($4.293$ lượt ghép dòng MinerU thành công trên $16.515$ đoạn văn bản). Song song với đó, cơ chế chuẩn hoá ký hiệu công thức ở cả phía truy vấn và chỉ mục từ khóa thưa đã nâng tỷ lệ tìm kiếm thành công phân đoạn tài liệu đích từ $1/12$ lên $11/12$ trên tập mẫu kiểm thử.
    
    \item \textbf{Mục tiêu 2 về thiết lập cơ sở tri thức véc-tơ toàn diện:} Hoàn thành xuất sắc. Toàn bộ chương trình môn Khoa học tự nhiên từ lớp 6 đến lớp 9 của cả ba bộ sách giáo khoa chuẩn mực hiện hành, bao gồm 12 cuốn sách tích hợp đầy đủ ba phân môn Vật lí, Hoá học và Sinh học, đã được số hoá và lập chỉ mục trọn vẹn: \textbf{2.387/2.399 trang tài liệu} nội dung bài học được xử lý thành \textbf{20.396 véc-tơ} ($16.515$ véc-tơ văn bản kích thước $1024$ chiều và $3.881$ véc-tơ hình ảnh kích thước $512$ chiều) lưu trữ trong ChromaDB, đồng bộ cùng chỉ mục thưa Okapi BM25 với $20.914$ từ vựng. Toàn bộ quy trình trích xuất, chuyển đổi và nạp dữ liệu này vận hành hiệu quả trực tiếp trên CPU của máy trạm thông thường mà không đòi hỏi tài nguyên GPU.

    \item \textbf{Mục tiêu 3 về tối ưu hóa pha truy xuất tri thức:} Hoàn thành xuất sắc và tạo nên đóng góp phương pháp luận cốt lõi. Kiểm nghiệm định lượng trên bộ dữ liệu kiểm thử $240$ câu hỏi lấy mẫu ngẫu nhiên ở quy mô vận hành thực tế khẳng định cấu hình đề xuất, gồm truy xuất lai kết hợp tái xếp hạng bằng mô hình tương tác chéo, đạt chỉ số \textbf{MRR đạt $0{,}7789$}, vượt trội so với các cấu hình đối chứng: chỉ mục thưa BM25 thuần đạt $0{,}6636$, kênh ngữ nghĩa thuần đạt $0{,}5664$ và phương pháp lai khi chưa tái xếp hạng đạt $0{,}6947$. Các thực nghiệm bóc tách thành phần cũng khẳng định vai trò tích cực của bước tái xếp hạng, đặc biệt nâng chỉ số R@1 của kênh véc-tơ đặc thêm $+53{,}6\%$, đồng thời xác nhận tác động bất lợi của cổng lọc khoảng cách tương đối ở quy mô thực tế, cung cấp cơ sở thực nghiệm vững chắc cho việc thiết lập cấu hình mặc định của hệ thống.

    \item \textbf{Mục tiêu 4 về thiết lập khung đánh giá thực nghiệm đối chiếu:} Hoàn thành toàn diện. Đề tài đã xây dựng thành công bộ dữ liệu kiểm thử $240$ câu hỏi có gắn nhãn nguồn xác thực và tiến hành đánh giá thực nghiệm đa chiều:
    \begin{itemize}
        \item Đối chiếu chi tiết giữa truy xuất lai và BM25 thuần cùng các cấu hình thành phần;
        \item Thực nghiệm đối sánh chuyên biệt giữa RAG đa phương thức và RAG thuần văn bản trên $52$ câu hỏi hình ảnh, mang lại kết luận định lượng quan trọng về việc duy trì cấu hình tối ưu;
        \item Kiểm chứng định lượng năng lực từ chối trả lời chính xác đối với các câu hỏi ngoài phạm vi môn học ($29/30$ câu, đạt tỷ lệ $96{,}67\%$, với điểm số giám khảo về Tính đúng đạt $4{,}87/5$), minh chứng cho khả năng kiểm soát ảo giác hiệu quả của hệ thống.
    \end{itemize}

    \item \textbf{Mục tiêu 5 về hiện thực hóa ứng dụng trợ lý học tập trực quan:} Hoàn thành đầy đủ. Ứng dụng web tương tác trên nền tảng Next.js đã được phát triển hoàn thiện, hỗ trợ người học gửi câu hỏi và nhận câu trả lời dưới dạng dòng dữ liệu phản hồi trực tiếp liên tục, hiển thị đồng bộ hình ảnh minh họa từ máy chủ, kết xuất \textbf{công thức khoa học chuẩn ký hiệu KaTeX} và cung cấp \textbf{khối trích dẫn minh bạch} phân tách rõ tên sách cùng số trang in thực tế.
\end{enumerate}

Về chất lượng câu trả lời sinh ra, đánh giá độc lập theo phương pháp mô hình ngôn ngữ lớn làm giám khảo với hội đồng bốn mô hình luân phiên ghi nhận điểm số trung bình toàn hệ thống đạt mức cao: Tính đúng đạt $4{,}033/5$, Độ trung thực đạt $4{,}417/5$ và Độ liên quan đạt $4{,}492/5$. Trong đó, nhóm câu hỏi ngoài phạm vi đạt điểm số tuyệt đối $5{,}00/5$ về Độ trung thực và Độ liên quan cùng điểm Tính đúng $4{,}87/5$; nhóm câu hỏi văn bản đạt Tính đúng $4{,}19/5$, Độ trung thực $4{,}56/5$ và Độ liên quan $4{,}59/5$. Đặc biệt, cơ chế xây dựng trích dẫn xác thực dựa trên siêu dữ liệu phân đoạn văn bản truy xuất đã triệt tiêu hoàn toàn hiện tượng mô hình ngôn ngữ tự tạo dựng số trang sai lệch, đảm bảo độ tin cậy khoa học vững chắc cho người học và giáo viên khi tra cứu học liệu.

\section{Hạn chế của đề tài}
Bên cạnh những kết quả đạt được, đề tài nhận diện các hạn chế kỹ thuật hiện tại cần được tiếp tục hoàn thiện (chi tiết số liệu tại Mục~\ref{sec:hanche}):
\begin{itemize}
    \item Quy trình kiểm định bộ dữ liệu kiểm thử $240$ câu hỏi hiện dựa trên một lượt rà soát chuyên gia tổng thể; việc bổ sung thêm quy trình thẩm định chéo độc lập từ nhiều chuyên gia chuyên ngành sẽ nâng cao hơn nữa độ chuẩn xác thống kê.
    \item Việc luân chuyển giữa bốn mô hình trong hội đồng giám khảo độc lập có thể tạo ra một biên độ dao động nhỏ về mặt phân bố điểm số.
    \item Quá trình ghép dòng công thức theo kiến trúc lai đã hoàn tất xử lý trên toàn bộ kho tri thức, song việc tiếp tục tiến hành các đợt đo tỷ lệ lỗi ký tự trên tập mẫu mở rộng sẽ giúp lượng hóa sâu sắc hơn mức độ chuẩn hóa của chỉ số dưới.
    \item Trường siêu dữ liệu số hiệu bài học hiện mới bao phủ $29{,}8\%$ số phân đoạn do hệ thống tuân thủ nguyên tắc tránh suy đoán khi mục lục bị gián đoạn bố cục.
    \item Hiện tượng suy giảm nhẹ tại chỉ số R@10 của cấu hình tái xếp hạng cần được tiếp tục phân tích qua các thực nghiệm quét ngưỡng điểm số tối thiểu.
    \item Cấu trúc hàng và cột của một số bảng biểu số liệu phức tạp đôi khi bị xáo trộn qua bước phân mảnh văn bản thông thường.
    \item Hệ thống mới được thử nghiệm chủ đạo qua các kịch bản kiểm thử định lượng tự động; quy mô triển khai thử nghiệm thực tế với giáo viên và học sinh tại các trường phổ thông còn giới hạn.
\end{itemize}

\section{Hướng phát triển tương lai}
Dựa trên những phát hiện khoa học và các mặt hạn chế được chỉ ra, các hướng nghiên cứu và phát triển tiếp theo của đề tài được định hình theo thứ tự ưu tiên như sau:

\begin{enumerate}
    \item \textbf{Triển khai thử nghiệm thực tế trong môi trường sư phạm:} Đưa hệ thống vào thử nghiệm diện rộng tại các cơ sở giáo dục Trung học cơ sở nhằm thu thập phản hồi định tính và định lượng từ đội ngũ giáo viên giảng dạy và học sinh, đánh giá tác động thực tiễn của trợ lý ảo đối với hiệu quả dạy và học môn Khoa học tự nhiên.
    \item \textbf{Đo lường và lượng hoá chuyên sâu chất lượng nhận dạng công thức:} Tiến hành các đợt thực nghiệm đối sánh mở rộng trên tập mẫu lớn văn bản sau ghép dòng để đánh giá định lượng mức cải thiện của chỉ số dưới và độ chuẩn xác của các phương trình hoá học, vật lí.
    \item \textbf{Nghiên cứu tối ưu hóa ngưỡng tham số của bộ tái xếp hạng:} Tiến hành quét lưới chi tiết đối với tham số \texttt{RERANK\_SCORE\_MIN} nhằm giải quyết triệt để hiện tượng nghịch đảo độ phủ tại chỉ số R@10, cân bằng tối ưu giữa độ chính xác và độ phủ của các phân đoạn tài liệu trích xuất.
    \item \textbf{Hoàn thiện trục phân tích mục lục bài học:} Nâng cấp bộ phân tích cấu trúc mục lục cho các bộ sách còn lại để hoàn thiện trường siêu dữ liệu số hiệu bài học đạt độ phủ $100\%$ trên toàn bộ kho tri thức, mở khoá tính năng tra cứu học liệu chuyên biệt theo từng bài học cụ thể.
    \item \textbf{Tích hợp kỹ thuật bảo toàn cấu trúc bảng biểu:} Ứng dụng các giải pháp nhận dạng và phân tích cấu trúc bảng biểu chuyên sâu nhằm bảo toàn nguyên vẹn mối quan hệ hàng và cột của dữ liệu số liệu khi phân mảnh văn bản.
    \item \textbf{Tích hợp mô hình thị giác -- ngôn ngữ trực tiếp:} Nghiên cứu kết hợp các mô hình thị giác -- ngôn ngữ đa phương thức để tiếp nhận trực tiếp hình ảnh sơ đồ và thí nghiệm phức tạp trong câu lệnh nhắc, khắc phục giới hạn của việc chỉ biểu diễn hình ảnh dưới dạng chuỗi siêu dữ liệu văn bản.
    \item \textbf{Lượng tử hóa và tối ưu hóa suy luận mô hình sinh:} Nghiên cứu ứng dụng các kỹ thuật lượng tử hóa trọng số nhằm triển khai mô hình ngôn ngữ lớn cục bộ với chi phí phần cứng tối thiểu, nâng cao tốc độ phản hồi và bảo mật dữ liệu học tập.
\end{enumerate}
